%%%%%%%%%%%%%%%%%%%%%%%%%%%%%%%%%%%%%%%%%%%%%%%%%%%%%%%%%%%%%%%%%%%%%%
% Lớp 12 | Ứng dụng đạo hàm để khảo sát hàm số | Khảo sát sự biến thiên và vẽ đồ thị hàm số
% Dạng 2D1?5-2: Các phép biến đổi đồ thị
% Số câu: 2  (sinh tự động bởi tach_cau.py; từ nay thầy sửa trực tiếp file này)
%%%%%%%%%%%%%%%%%%%%%%%%%%%%%%%%%%%%%%%%%%%%%%%%%%%%%%%%%%%%%%%%%%%%%%

%%%%%%%%%%%%%%%% TRẮC NGHIỆM 4 PHƯƠNG ÁN %%%%%%%%%%%%%%%%

% nguồn: Phan1_Lop_12_Dai_so_Chuong_1
\begin{ex}%[2D1H5-2]%[ma:c8906312]
		Cho hàm số $y=f(x)$ xác định, liên  tục trên $\mathbb{R}$ và có đồ thị như hình bên.
		\begin{center}
			\begin{tikzpicture}[line join=round,line cap=round,>=stealth]
				\tikzset{ethicknode/.style={scale=0.9}} 
				\draw[->] (-2.5,0)--(2,0) node[below] {$x$}; 
				 \draw[->] (0,-1.1)--(0,3.1) node[below left] {$y$};
				\draw (0,0) node [below left] {$O$}; 
				\foreach \x/\nx in  {1/1}
				\draw[thin] (\x,1pt)--(\x,-1pt) node  [below] {$\nx$};
				\foreach \y/\ny  in {1/1}
				\draw[thin] (1pt,\y)--(-1pt,\y) node  [left] {$\ny$};
				\begin{scope}
					\clip (-2.5,-1) rectangle (1.5,3); 
					 \draw[samples=200,domain=-2.5:2.5,smooth,variable=\x] plot (\x,{3^(\x)});
				\end{scope}
			\end{tikzpicture}
		\end{center}
		Đồ thị nào dưới đây là đồ thị của  hàm số $y=f(x)+1$?
	\choice 
			{\begin{tikzpicture}[line join=round,line cap=round,>=stealth]
		 \tikzset{ethicknode/.style={scale=0.9}}
		\draw[->] (-2.5,0)--(2,0) node[below] {$x$}; 
		\draw[->]  (0,-1.1)--(0,3.1) node[below left] {$y$};
		\draw (0,0) node [below left] {$O$}; 
		 \foreach \x/\nx in {1/1}
		\draw[thin] (\x,1pt)--(\x,-1pt) node [below] {$\nx$}; 
		\foreach \y/\ny in {1/1} 
		 \draw[thin] (1pt,\y)--(-1pt,\y) node [left] {$\ny$};
		\begin{scope}
			\clip  (-2.5,-1) rectangle (1.5,3);
			\draw[samples=200,domain=-2.5:2.5,smooth,variable=\x] plot (\x,{3^(\x)}); 
		\end{scope}
	\end{tikzpicture}
			} 
			{\True \begin{tikzpicture}[line join=round,line cap=round,>=stealth]
		 \tikzset{ethicknode/.style={scale=0.9}}
		 \draw[->] (-2.5,0)--(2,0) node[below] {$x$};
		\draw[->] (0,-1.1)--(0,3.1)  node[below left] {$y$};
		\draw (0,0) node [below left] {$O$}; 
		 \foreach \x/\nx in {1/1}
		\draw[thin] (\x,1pt)--(\x,-1pt)  node [below] {$\nx$};
		\foreach \y/\ny in  {1/1,2/2}
		\draw[thin] (1pt,\y)--(-1pt,\y) node  [left] {$\ny$};
		\begin{scope}
			\clip  (-2.5,-1) rectangle (1.5,3);
			\draw[samples=200,domain=-2.5:2.5,smooth,variable=\x]  plot (\x,{3^(\x)+1});
		\end{scope}
	\end{tikzpicture}
			 }
			{\begin{tikzpicture}[line join=round,line cap=round,>=stealth]
		 \tikzset{ethicknode/.style={scale=0.9}}
		 \draw[->] (-2.5,0)--(2,0) node[below] {$x$};
		 \draw[->] (0,-1.1)--(0,3.1) node[below left] {$y$};
		\draw (0,0) node  [below right] {$O$};
		\foreach  \x/\nx in {1/1}
		\draw[thin] (\x,1pt)--(\x,-1pt)  node [below] {$\nx$};
		\foreach  \y/\ny in {1/1}
		\draw[thin] (1pt,\y)--(-1pt,\y) node [left]  {$\ny$};
		\begin{scope}
			\clip (-2.5,-1) rectangle  (1.5,3);
			\draw[samples=200,domain=-2.5:2.5,smooth,variable=\x] plot (\x,{3^(\x)-1}); 
		\end{scope}
	\end{tikzpicture}
			 }
			{\begin{tikzpicture}[line join=round,line cap=round,>=stealth]
		\tikzset{ethicknode/.style={scale=0.9}} 
		\draw[->]  (-2.5,0)--(2,0) node[below] {$x$};
		\draw[->] (0,-1.1)--(0,3.1) node[below left]  {$y$};
		\draw (0,0) node  [below left] {$O$};
		 \foreach \x/\nx in {1/1}
		\draw[thin] (\x,1pt)--(\x,-1pt)  node [below] {$\nx$};
		\foreach \y/\ny  in {1/1,2/2}
		 \draw[thin] (1pt,\y)--(-1pt,\y) node [left] {$\ny$};
		\begin{scope}
			 \clip (-2.5,-1) rectangle (1.5,3);
			 \draw[samples=200,domain=-2.5:2.5,smooth,variable=\x] plot (\x,{3^(\x+0.7)});
		\end{scope}
	\end{tikzpicture}
			} 
	\loigiai{ 
			Từ đồ thị hàm  số $y=f(x)$ tìm đồ thị hàm số $y=f(x)+a$.\\
			Tịnh tiến đồ thị hàm  số $y=f(x)$ lên phía trên (theo phương $Oy$) $a$ đơn vị nếu $a>0$.\\
			Hàm số $y=f(x)+1 \Rightarrow a=1>0$ nên ta được đồ thị hàm số cần tìm tịnh tiến đồ  thị hàm số $y=f(x)$ lên phía trên (theo phương $Oy$) $1$ đơn vị.
		} 
	\end{ex}


%%%%%%%%%%%%%%%% ĐÚNG/SAI %%%%%%%%%%%%%%%%

% nguồn: Phan1_Lop_12_Dai_so_Chuong_1
% nguồn: 2026-11-16_12_2D1_TL_Ham_so (nhập 16/11/2026)
\begin{bt}%[2D1H5-2]%[ma:f2951780]
% nguon: Archive/2024/TaiLieuGiangDay/1.HamSo/1. K12 - 4in1 - Hamso/ThubaiGoc/Bai5-Dang4.tex #2234
[Mức độ 1]%[Bài giảng 12 mới, Lê Hồng Phi]
Đồ thị hàm số $y=\log_2 x$ và đường thẳng $x=-1$ cắt nhau tại bao nhiêu điểm?
\loigiai
{Đồ thị hàm số $y=\log_2 x$  nằm phía bên phải trục $Oy$.\\
Đường thẳng $x=-1$ nằm bên trái trục $Oy$ nên đồ thị hàm số đã cho và đường thẳng $x=-1$ không cắt nhau.
}
\end{bt}

% nguồn: 2026-11-16_12_2D1_TL_Ham_so (nhập 16/11/2026)
% nguồn: 2026-11-16_12_2D1_TL_Ham_so (nhập 16/11/2026)
\begin{bt}%[2D1H5-2]%[ma:143e7e23]
% nguon: Archive/2024/TaiLieuGiangDay/1.HamSo/1. K12 - 4in1 - Hamso/ThubaiGoc/Bai5-Dang5.tex #2305

	Cho đồ thị hàm số $y=x^3-6x^2+9x$ như hình dưới. Hãy suy ra đồ thị của hàm số $y=|x|^3-6x^2+9|x|$.
	
	\begin{center}
		\begin{tikzpicture}[>=stealth,x=1cm,y=1cm,scale=0.4]
		\def\a{1} % Hệ số a phải khác 0
		\def\b{-6}
		\def\c{9}
		\def\d{0}
		\draw[->] (-5,0) -- (5,0) node[below]{\scriptsize $x$};
		\foreach \x in {-4,-3,-2,-1,1,2,3,4}
		\draw[shift={(\x,0)}] (0,2pt)--(0,-2pt) node[below]{};
		\draw[->] (0,-1) -- (0,5)node[right]{\scriptsize $y$};
		\foreach \y in {1,2,3,4,5}
		\draw[shift={(0,\y)}] (2pt,0)--(-2pt,0) node[right]{};
		\fill (0,0)node[below left]{\scriptsize $O$} circle(2pt);
		\clip (-5,-2.5)rectangle(5,5);
		\draw[thick,samples=150,smooth,domain=-0.1:5] plot(\x,{\a*(\x)^3+(\b)*(\x)^2+(\c)*\x+(\d)});
		\draw (0,-2) node{$\text{Hình } 1$};
		\end{tikzpicture}
	\end{center}
	\loigiai{
		Tập xác định $\mathscr{D}=\mathbb{R}$.\\
		Ta thấy đồ thị hàm số $y=|x|^3-6|x|^2+9|x|$ có được bằng cách:
		\begin{itemize}
			\item Giữ nguyên phần đồ thị của hàm số $y=x^3-6x^2+9x$ nằm bên phải trục $Oy$.
			\item Lấy đối xứng qua trục tung phần đồ thị của $y=x^3-6x^2+9x$ nằm bên phải trục $Oy$.
		\end{itemize}
		Do đó đồ thị hàm số $y=|x|^3-6|x|^2+9|x|$ là
		\begin{center}
			
			\begin{tikzpicture}[>=stealth,x=1cm,y=1cm,scale=0.4]
				\def\a{1} % Hệ số a phải khác 0
				\def\b{-6}
				\def\c{9}
				\def\d{0}
				\draw[->] (-5,0) -- (5,0) node[below]{\scriptsize $x$};
				\foreach \x in {-4,-3,-2,-1,1,2,3,4}
				\draw[shift={(\x,0)}] (0,2pt)--(0,-2pt) node[below]{};
				\draw[->] (0,-1) -- (0,5)node[right]{\scriptsize $y$};
				\foreach \y in {1,2,3,4,5}
				\draw[shift={(0,\y)}] (2pt,0)--(-2pt,0) node[right]{};
				\fill (0,0)node[below left]{\scriptsize $O$} circle(2pt);
				\clip (-5,-2.5)rectangle(5,5);
				\draw[thick,samples=250,smooth,domain=-5:5] plot(\x,{\a*abs((\x))^3+(\b)*(\x)^2+(\c)*abs(\x)+(\d)});
				\draw (0,-2) node{$\text{Hình } 2$};
			\end{tikzpicture}			
		\end{center}
	}

\end{bt}

% nguồn: 2026-11-16_12_2D1_TL_Ham_so (nhập 16/11/2026)
\begin{bt}%[2D1H5-2]%[ma:4a8b71d8]
% nguon: Archive/2024/TaiLieuGiangDay/1.HamSo/1. K12 - 4in1 - Hamso/ThubaiGoc/Bai5-Dang5.tex #2306
 
	Cho hàm số $y=f(x)=\dfrac{x^2-2x+1}{x+1}$ có đồ thị như hình dưới. Hãy suy ra đồ thị của hàm số $y=f(|x|)$?
	\begin{center}
			\begin{tikzpicture}[>=stealth,line join=round,line cap=round,font=\footnotesize,scale=.4]
			\tikzset{declare function={
					x1=-6.5;x2=4.5;y1=-12.5;y2=5.5;% Miền hình vẽ
					a=1;b=-2;c=1; m=1;n=1;
					f(\x)=(a*(\x)^2+b*(\x)+c)/(m*(\x)+n);	
					a1=a/m;b1=(b-a*n/m)/m;
					g(\x)=a1*\x+b1;
					q1=-(n/m)-.1;q2=-n/m+.1;
			}}
			\begin{scope}
				\clip (x1,y1)rectangle(x2,y2);
				\draw [smooth,samples=300,domain=x1:q1]plot(\x,{f(\x)});
				\draw [smooth,samples=300,domain=q2:x2]plot(\x,{f(\x)});
				\draw [smooth,samples=100,domain=x1:x2,dashed]plot(\x,{g(\x)});
				\draw[dashed] (-n/m,y1)--(-n/m,y2);
			\end{scope}
			\draw[->] (0,y1)--(0,y2)node[left]{$y$};
			\draw[->] (x1,0)--(x2,0)node[above]{$x$};
			\filldraw[black](0,0)circle(1.2pt)node[above left]{$O$};
			\draw[dashed]
			(-3,0)|-(0,-8)
			;
			
			\foreach \x/\goc in {-3/-90,-1/-90,1/-70,2/-90}{
				\draw[fill=black] (\x,0)circle(1.2pt) node[shift={(\goc:2.8mm)}]{$\x$};
			}
			\foreach \y/\goc in {-3,-4,1,2,-8}{
				\draw[fill=black] (0,\y)circle(1.2pt)node[shift={(\goc:2.8mm)}]{$\y$};
			}
		\end{tikzpicture}
	\end{center}
	\loigiai{
		Tập xác định $\mathscr{D}=\mathbb{R}\setminus {-1}$.\\
		Ta thấy đồ thị hàm số $y=f(|x|)$ có được bằng cách:
		\begin{itemize}
			\item Giữ nguyên phần đồ thị của hàm số $y=f(x)=\dfrac{x^2-2x+1}{x+1}$ nằm bên phải trục $Oy$.
			\item Lấy đối xứng qua trục tung phần đồ thị của $y=f(x)=\dfrac{x^2-2x+1}{x+1}$ nằm bên phải trục $Oy$.
		\end{itemize}
		Do đó đồ thị hàm số $y=f(|x|)=\dfrac{|x|^2-2|x|+1}{|x|+1}$ là
		\begin{center}
		\begin{tikzpicture}[>=stealth,line join=round,line cap=round,font=\footnotesize,scale=.4]
			\tikzset{declare function={
					x1=-6.5;x2=4.5;y1=-4.5;y2=5.5;% Miền hình vẽ
					a=1;b=-2;c=1; m=1;n=1;
					f(\x)=(a*(\x)^2+b*(\x)+c)/(m*(\x)+n);	
					a1=a/m;b1=(b-a*n/m)/m;
					g(\x)=a1*\x+b1;
					q1=-(n/m)-.1;q2=-n/m+.1;
			}}
			\begin{scope}
				\clip (x1,y1)rectangle(x2,y2);
				\draw [smooth,samples=300,domain=x1:q1]plot(\x,{f(abs(\x))});
				\draw [smooth,samples=300,domain=q2:x2]plot(\x,{f(abs(\x))});
				
				\draw [smooth,samples=100,domain=x1:x2,dashed]plot(\x,{g(abs(\x))});
				%		\draw[dashed] (-n/m,y1)--(-n/m,y2);
			\end{scope}
			\draw[->] (0,y1)--(0,y2)node[left]{$y$};
			\draw[->] (x1,0)--(x2,0)node[above]{$x$};
			\filldraw[black](0,0)circle(1.2pt)node[below right]{$O$};
			\draw[dashed]
			%	(-3,0)|-(0,-8)
			;
			
			\foreach \x/\goc in {-3/-90,-1/-90,1/-70,2/-90}{
				\draw[fill=black] (\x,0)circle(1.2pt) node[shift={(\goc:2.8mm)}]{$\x$};
			}
			\foreach \y/\goc in {1,2,-3}{
				\draw[fill=black] (0,\y)circle(1.2pt)node[shift={(\goc:2.8mm)}]{$\y$};
			}
		\end{tikzpicture}
	\end{center}
	}

\end{bt}

% nguồn: 2026-11-16_12_2D1_TL_Ham_so (nhập 16/11/2026)
\begin{bt}%[2D1H5-2]%[ma:6d41460b]
% nguon: Archive/2024/TaiLieuGiangDay/1.HamSo/1. K12 - 4in1 - Hamso/ThubaiGoc/Bai5-Dang5.tex #2307

	Cho hàm số $y=log_4 x$ có đồ thị như hình sau. Hãy suy ra đồ thị của hàm số $y=|log_4 x|$.
	\begin{center}
		\begin{tikzpicture}[>=stealth,line join=round,line cap=round,font=\footnotesize,scale=.8]
			\tikzset{declare function={
					x1=-1.5;x2=3.5;y1=-2.5;y2=2.5;% Miền hình vẽ
					a=4;b=-2;c=1; m=1;n=1;
					f(\x)=ln(\x)/ln(a);	
					x0=0.02;
			}}
			\begin{scope}
				\clip (x1,y1)rectangle(x2,y2);
				\draw [smooth,samples=300,domain=x0:x2]plot(\x,{f(\x)});
			\end{scope}
			\draw[->] (0,y1)--(0,y2)node[left]{$y$};
			\draw[->] (x1,0)--(x2,0)node[above]{$x$};
			\filldraw[black](0,0)circle(1.2pt)node[above left]{$O$};
			\draw[dashed];
			
			\foreach \x/\goc in {-1/-90,1/-70,2/-90}{
				\draw[fill=black] (\x,0)circle(1.2pt) node[shift={(\goc:2.8mm)}]{$\x$};
			}
			\foreach \y/\goc in {1/180}{
				\draw[fill=black] (0,\y)circle(1.2pt)node[shift={(\goc:2.8mm)}]{$\y$};
			}
		\end{tikzpicture}
		
	\end{center}
	\loigiai{
	Tập xác định $\mathscr{D}=(0;+\infty)$.\\
	Đồ thị hàm số $y=|log_4 x|$ có được bằng cách:
	\begin{itemize}
		\item Giữ nguyên phần đồ thị của hàm số $y=log_4 x$ nằm trên trục hoành.
		\item Lấy đối xứng qua trục hoành phần đồ thị của $y=log_4 x$ nằm dưới trục hoành.
	\end{itemize}
	Do đó đồ thị của hàm số $y=|log_4 x|$ là	
	\begin{center}
		\begin{tikzpicture}[>=stealth,line join=round,line cap=round,font=\footnotesize,scale=.8]
			\tikzset{declare function={
					x1=-1.5;x2=3.5;y1=-2.5;y2=2.5;% Miền hình vẽ
					a=4;b=-2;c=1; m=1;n=1;
					f(\x)=ln(\x)/ln(a);	
					x0=0.02;
			}}
			\begin{scope}
				\clip (x1,y1)rectangle(x2,y2);
				\draw [smooth,samples=300,domain=x0:x2]plot(\x,{abs(f(\x))});
			\end{scope}
			\draw[->] (0,y1)--(0,y2)node[left]{$y$};
			\draw[->] (x1,0)--(x2,0)node[above]{$x$};
			\filldraw[black](0,0)circle(1.2pt)node[above left]{$O$};
			\draw[dashed];
			
			\foreach \x/\goc in {-1/-90,1/-70,2/-90}{
				\draw[fill=black] (\x,0)circle(1.2pt) node[shift={(\goc:2.8mm)}]{$\x$};
			}
			\foreach \y/\goc in {1/180}{
				\draw[fill=black] (0,\y)circle(1.2pt)node[shift={(\goc:2.8mm)}]{$\y$};
			}
		\end{tikzpicture}
	\end{center}
	}
\end{bt}

% nguồn: 2026-11-16_12_2D1_TL_Ham_so (nhập 16/11/2026)
\begin{bt}%[2D1H5-2]%[ma:6554b8bb]
% nguon: Archive/2024/TaiLieuGiangDay/1.HamSo/1. K12 - 4in1 - Hamso/ThubaiGoc/Bai5-Dang5.tex #2308
 
	Cho hàm số $y=\left(\dfrac{1}{3}\right)^x$ có đồ thị như hình dưới. Hãy suy ra đồ thị của hàm số $y=\left(\dfrac{1}{3}\right)^{|x|}$?
	\begin{center}
		\begin{tikzpicture}[>=stealth,line join=round,line cap=round,font=\footnotesize,scale=.8]
			\tikzset{declare function={
					x1=-2.5;x2=3.5;y1=-1.5;y2=2.5;% Miền hình vẽ
					a=1/3;
					f(\x)=a^(\x);	
			}}
			\begin{scope}
				\clip (x1,y1)rectangle(x2,y2);
				\draw [smooth,samples=300,domain=x1:x2]plot(\x,{f(\x)});
			\end{scope}
			\draw[->] (0,y1)--(0,y2)node[left]{$y$};
			\draw[->] (x1,0)--(x2,0)node[above]{$x$};
			\filldraw[black](0,0)circle(1.2pt)node[above left]{$O$};
			\draw[dashed];
			
			\foreach \x/\goc in {-2/-90,-1/-90,1/-70,2/-90}{
				\draw[fill=black] (\x,0)circle(1.2pt) node[shift={(\goc:2.8mm)}]{$\x$};
			}
			\foreach \y/\goc in {1/180}{
				\draw[fill=black] (0,\y)circle(1.2pt)node[shift={(\goc:2.8mm)}]{$\y$};
			}
		\end{tikzpicture}
	\end{center}
	\loigiai{
		Tập xác định $\mathscr{D}=\mathbb{R}$.\\
		Ta thấy đồ thị hàm số $y=\left(\dfrac{1}{3}\right)^{|x|}$ có được bằng cách:
		\begin{itemize}
			\item Giữ nguyên phần đồ thị của hàm số $y=\left(\dfrac{1}{3}\right)^x$ nằm bên phải trục $Oy$.
			\item Lấy đối xứng qua trục tung phần đồ thị của $y=\left(\dfrac{1}{3}\right)^x$ nằm bên phải trục $Oy$.
		\end{itemize}
		Do đó đồ thị hàm số $y=\left(\dfrac{1}{3}\right)^{|x|}$ là
		\begin{center}
		\begin{tikzpicture}[>=stealth,line join=round,line cap=round,font=\footnotesize,scale=.8]
			\tikzset{declare function={
					x1=-2.5;x2=3.5;y1=-1.5;y2=2.5;% Miền hình vẽ
					a=1/3;
					f(\x)=a^(\x);	
			}}
			\begin{scope}
				\clip (x1,y1)rectangle(x2,y2);
				\draw [smooth,samples=300,domain=x1:x2]plot(\x,{f(abs(\x))});
			\end{scope}
			\draw[->] (0,y1)--(0,y2)node[left]{$y$};
			\draw[->] (x1,0)--(x2,0)node[above]{$x$};
			\filldraw[black](0,0)circle(1.2pt)node[above left]{$O$};
			\draw[dashed];
			
			\foreach \x/\goc in {-2/-90,-1/-90,1/-70,2/-90}{
				\draw[fill=black] (\x,0)circle(1.2pt) node[shift={(\goc:2.8mm)}]{$\x$};
			}
			\foreach \y/\goc in {1/180}{
				\draw[fill=black] (0,\y)circle(1.2pt)node[shift={(\goc:2.8mm)}]{$\y$};
			}
		\end{tikzpicture}
		\end{center}
	}
	
\end{bt}

% nguồn: 2026-11-16_12_2D1_TL_Ham_so (nhập 16/11/2026)
\begin{bt}%[2D1V5-2]%[ma:7f871b5f]
% nguon: Archive/2024/TaiLieuGiangDay/1.HamSo/1. K12 - 4in1 - Hamso/ThubaiGoc/Bai5-Dang5.tex #2309
 
	Cho hàm số $y=-x^3+3x^2-2$ có đồ thị như hình sau. Từ đó suy ra đồ thị của hàm số $y=|-|x|^3+3|x|^2-2|$ ?
	\begin{center}
		\begin{tikzpicture}[>=stealth,line join=round,line cap=round,font=\footnotesize,scale=.6]
			\tikzset{declare function={
					x1=-4.5;x2=4.5;y1=-2.5;y2=3.5;
					a=-1;b=3;c=0;d=-2;
					f(\x)=a*(\x)^3+b*(\x)^2+c*(\x)+d;
			}}
			\begin{scope}
				\clip (x1,y1)rectangle(x2,y2);
				\draw [smooth,samples=100,domain=x1:x2]plot(\x,{f(\x)});
			\end{scope}
			\draw[->] (0,y1)--(0,y2)node[left]{$y$};
			\draw[->] (x1,0)--(x2,0)node[above]{$x$};
			\filldraw(0,0)circle(1.2pt)node[below left]{$O$};
			\draw[dashed] 
			%	(1,0)|-(0,-2) (3,0)|-(0,-4)
			;
			\foreach \x/\goc in {-2/-110,-3/-110,-1/-110,1/-70,2/-90,3/-90}{
				\draw[fill=black] (\x,0)circle(1.2pt) node[shift={(\goc:2.8mm)}]{$\x$};
			}
			\foreach \y/\goc in {-2/180,1/180,2/180,3/180}{
				\draw[fill=black] (0,\y)circle(1.2pt)node[shift={(\goc:2.8mm)}]{$\y$};
			}
		\end{tikzpicture}
	\end{center}
	\loigiai{
		Tập xác định $\mathscr{D}=\mathbb{R}$.\\
		Ta thấy đồ thị hàm số $y=|-|x|^3+3|x|^2-2|$ có được bằng cách:
		\begin{itemize}
			\item Giữ nguyên phần đồ thị của hàm số $y=-x^3+3x^2-2$ nằm bên phải trục $Oy$.
			\item Lấy đối xứng qua trục tung phần đồ thị của $y=-x^3+3x^2-2$ nằm bên phải trục $Oy$.
			\item Sau đó, giữ nguyên phần đồ thị nằm phía trên trục hoành.
			\item Lấy đối xứng qua trục hoành phần đồ thị nằm bên dưới trục hoành.
		\end{itemize}
		Do đó đồ thị hàm $y=|-|x|^3+3|x|^2-2|$ là 
		\begin{center}
			\begin{tikzpicture}[>=stealth,line join=round,line cap=round,font=\footnotesize,scale=.6]
				\tikzset{declare function={
						x1=-4.5;x2=4.5;y1=-2.5;y2=3.5;
						a=-1;b=3;c=0;d=-2;
						f(\x)=a*(\x)^3+b*(\x)^2+c*(\x)+d;
				}}
				\begin{scope}
					\clip (x1,y1)rectangle(x2,y2);
					\draw [smooth,samples=500,domain=x1:x2]plot(\x,{abs(f(abs(\x)))});
				\end{scope}
				\draw[->] (0,y1)--(0,y2)node[left]{$y$};
				\draw[->] (x1,0)--(x2,0)node[above]{$x$};
				\filldraw(0,0)circle(1.2pt)node[below left]{$O$};
				\draw[dashed] 
				%	(1,0)|-(0,-2) (3,0)|-(0,-4)
				;
				\foreach \x/\goc in {-2/-110,-3/-110,-1/-110,1/-70,2/-90,3/-90}{
					\draw[fill=black] (\x,0)circle(1.2pt) node[shift={(\goc:2.8mm)}]{$\x$};
				}
				\foreach \y/\goc in {-2/180,1/180,2/180,3/180}{
					\draw[fill=black] (0,\y)circle(1.2pt)node[shift={(\goc:2.8mm)}]{$\y$};
				}
			\end{tikzpicture}
		\end{center}
	}	
\end{bt}

% nguồn: 2026-11-16_12_2D1_TL_Ham_so (nhập 16/11/2026)
\begin{bt}%[2D1V5-2]%[ma:00ee544c]
% nguon: Archive/2024/TaiLieuGiangDay/1.HamSo/1. K12 - 4in1 - Hamso/ThubaiGoc/Bai5-Dang5.tex #2310

	Cho hàm số $y=\dfrac{x^2+x-4}{x-1}$ có đồ thị như hình sau. Từ đó suy ra đồ thị của hàm số $y=\left|\dfrac{|x|^2+|x|-4}{|x|-1}\right|$ ?
	\begin{center}
	\begin{tikzpicture}[>=stealth,line join=round,line cap=round,font=\footnotesize,scale=.6]
		\tikzset{declare function={
				x1=-3.5;x2=4.5;y1=-3.5;y2=6.5;
				a=1;b=1;c=-4;m=1;n=-1;
				f(\x)=(a*(\x)^2+b*\x+c)/(m*(\x)+n);
				a1=a/m;b1=b-(a*n)/m*m;
				g(\x)=a1*\x+b1;
				x01=-n/m-.02;x02=-n/m+.02;
		}}
		\begin{scope}
			\clip (x1,y1)rectangle(x2,y2);
			\draw [smooth,samples=100,domain=x1:x01]plot({\x},{f(\x)});
			\draw [smooth,samples=100,domain=x02:x2]plot({\x},{f((\x))});
			\draw [smooth,samples=100,domain=x1:x2,dashed]plot({\x},{g((\x))});
			\draw[dashed] (-n/m,y1)--(-n/m,y2);
		\end{scope}
		\draw[->] (0,y1)--(0,y2)node[left]{$y$};
		\draw[->] (x1,0)--(x2,0)node[above]{$x$}; % Corrected label
		\filldraw(0,0)circle(1.2pt)node[above left]{$O$};
		\foreach \x/\goc in {1/-90,-1/-90}{
			\draw[fill=black] (\x,0)circle(1.2pt) node[shift={(\goc:2.8mm)}]{$\x$};
		}
		\foreach \y/\goc in {1,2,3,4}{
			\draw[fill=black] (0,\y)circle(1.2pt)node[shift={(\goc:2.8mm)}]{$\y$};
		}
	\end{tikzpicture}
	
	\end{center}
	\loigiai{
		Tập xác định $\mathscr{D}=\mathbb{R}\setminus {1}$.\\
		Ta thấy đồ thị hàm số $y=\left|\dfrac{|x|^2+|x|-4}{|x|-1}\right|$ có được bằng cách:
		\begin{itemize}
			\item Giữ nguyên phần đồ thị của hàm số $y=\dfrac{x^2+x-4}{x-1}$ nằm bên phải trục $Oy$.
			\item Lấy đối xứng qua trục tung phần đồ thị của $y=\dfrac{x^2+x-4}{x-1}$ nằm bên phải trục $Oy$.
			\item Sau đó, giữ nguyên phần đồ thị nằm phía trên trục hoành.
			\item Lấy đối xứng qua trục hoành phần đồ thị nằm bên dưới trục hoành.
		\end{itemize}
		Do đó đồ thị hàm $y=\left|\dfrac{|x|^2+|x|-4}{|x|-1}\right|$ là 
		\begin{center}
			\begin{tikzpicture}[>=stealth,line join=round,line cap=round,font=\footnotesize,scale=.6]
				\tikzset{declare function={
						x1=-4.5;x2=4.5;y1=-3.5;y2=6.5;
						a=1;b=1;c=-4;m=1;n=-1;
						f(\x)=(a*(\x)^2+b*\x+c)/(m*(\x)+n);
						a1=a/m;b1=b-(a*n)/m*m;
						g(\x)=a1*\x+b1;
						x01=-n/m-.02;x02=-n/m+.02;
				}}
				\begin{scope}
					\clip (x1,y1)rectangle(x2,y2);
					\draw [smooth,samples=160,domain=x1:x01]plot({\x},{abs(f(abs(\x)))});
					\draw [smooth,samples=200,domain=x02:x2]plot({\x},{abs(f(abs(\x)))});
					\draw [smooth,samples=100,domain=x1:x2,dashed]plot({\x},{g(\x)});
					\draw [smooth,samples=100,domain=x1:x2,dashed]plot({\x},{g(-\x)});
					\draw[dashed] (-n/m,y1)--(-n/m,y2);
					\draw[dashed] (n/m,y1)--(n/m,y2);
				\end{scope}
				\draw[->] (0,y1)--(0,y2)node[left]{$y$};
				\draw[->] (x1,0)--(x2,0)node[above]{$x$}; % Corrected label
				\filldraw(0,0)circle(1.2pt)node[above left]{$O$};
				\foreach \x/\goc in {1/-90,-1/-90}{
					\draw[fill=black] (\x,0)circle(1.2pt) node[shift={(\goc:2.8mm)}]{$\x$};
				}
				\foreach \y/\goc in {1,2,3,4}{
					\draw[fill=black] (0,\y)circle(1.2pt)node[shift={(\goc:2.8mm)}]{$\y$};
				}
			\end{tikzpicture}
			
		\end{center}	
	}	
\end{bt}

% nguồn: 2026-11-16_12_2D1_TL_Ham_so (nhập 16/11/2026)
\begin{bt}%[2D1N5-2]%[ma:7a3ea874]
% nguon: Archive/2024/TaiLieuGiangDay/1.HamSo/1. K12 - 4in1 - Hamso/ThubaiGoc/Bai5-Dang5.tex #2314

	\immini{Cho hàm số $y=f(x)$ có đồ thị như hình vẽ sau. Hãy xác định khoảng đồng biến của hàm số $y=f(|x|)$.}
	{
		\begin{tikzpicture}[>=stealth,line join=round,line cap=round,font=\footnotesize,scale=.6]
			\tikzset{declare function={
					x1=-3.5;x2=4.5;y1=-2.5;y2=4.5;% Miền hình vẽ
					a=1;b=-1;c=-1;d=1; m=1;n=1;
					f(\x)=a*(\x)^3+b*(\x)^2+c*(\x)+d;	
					a1=a/m;b1=(b-a*n/m)/m;
					g(\x)=a1*\x+b1;
					q1=-(n/m)-.1;q2=-n/m+.1;
			}}
			\begin{scope}
				\clip (x1,y1)rectangle(x2,y2);
				\draw [smooth,samples=100,domain=x1:x2]plot(\x,{f(\x)});
			\end{scope}
			\draw[->] (0,y1)--(0,y2)node[left]{$y$};
			\draw[->] (x1,0)--(x2,0)node[above]{$x$};
			\filldraw(0,0)circle(1.2pt)node[below left]{$O$};
			`	\draw[dashed]
			;
			% Đánh dấu tọa độ trên trục ox và oy
			\foreach \x/\goc in {-3/-90,-1/-90,1/-70,2/-90}{
				\draw[fill=black] (\x,0)circle(1.2pt) node[shift={(\goc:2.8mm)}]{$\x$};
			}
			\foreach \y/\goc in {-2,1,2}{
				\draw[fill=black] (0,\y)circle(1.2pt)node[shift={(\goc:2.8mm)}]{$\y$};
			}
		\end{tikzpicture}
	}
	\loigiai{	
		Ta thấy đồ thị hàm số $y=f(|x|)$ có được bằng cách:
		\begin{itemize}
			\item Giữ nguyên phần đồ thị của hàm số $y=f(x)$ nằm bên phải trục $Oy$.
			\item Lấy đối xứng qua trục tung phần đồ thị của $y=f(x)$ nằm bên phải trục $Oy$.
		\end{itemize}
		Do đó đồ thị hàm số $y=f(|x|)$ là
		\begin{center}
			\begin{tikzpicture}[>=stealth,line join=round,line cap=round,font=\footnotesize,scale=.6]
				\tikzset{declare function={
						x1=-3.5;x2=4.5;y1=-2.5;y2=4.5;% Miền hình vẽ
						a=1;b=-1;c=-1;d=1; m=1;n=1;
						f(\x)=a*(\x)^3+b*(\x)^2+c*(\x)+d;	
						a1=a/m;b1=(b-a*n/m)/m;
						g(\x)=a1*\x+b1;
						q1=-(n/m)-.1;q2=-n/m+.1;
				}}
				\begin{scope}
					\clip (x1,y1)rectangle(x2,y2);
					\draw [smooth,samples=100,domain=x1:x2]plot(\x,{f(abs(\x))});
				\end{scope}
				\draw[->] (0,y1)--(0,y2)node[left]{$y$};
				\draw[->] (x1,0)--(x2,0)node[above]{$x$};
				\filldraw(0,0)circle(1.2pt)node[below left]{$O$};
				`	\draw[dashed]
				;
				% Đánh dấu tọa độ trên trục ox và oy
				\foreach \x/\goc in {-3/-90,-1/-90,1/-70,2/-90}{
					\draw[fill=black] (\x,0)circle(1.2pt) node[shift={(\goc:2.8mm)}]{$\x$};
				}
				\foreach \y/\goc in {-2,1,2}{
					\draw[fill=black] (0,\y)circle(1.2pt)node[shift={(\goc:2.8mm)}]{$\y$};
				}
			\end{tikzpicture}
		\end{center}
		Từ đồ thị hàm số $y=f(\left|x\right|)$ ta thấy hàm số $y=f(\left|x\right|)$ đồng biến trên khoảng $(-1;0)$ và $(1;+\infty)$.	
		
	}
\end{bt}

% nguồn: 2026-11-16_12_2D1_TL_Ham_so (nhập 16/11/2026)
% nguồn: 2026-11-16_12_2D1_TL_Ham_so (nhập 16/11/2026)
\begin{bt}%[2D1H5-2]%[ma:be70f8a4]
% nguon: Archive/2024/TaiLieuGiangDay/1.HamSo/1. K12 - 4in1 - Hamso/ThubaiGoc/Bai5-Dang5.tex #2316

	\immini{
		Cho hàm số $ y=\dfrac{2x-1}{x-1}$ có đồ thị như hình vẽ. Tìm các khoảng đồng biến của hàm số $ y=\dfrac{2\left|x\right|-1}{\left|x\right|-1}$. 
	}
	{\begin{tikzpicture}[>=stealth,line join=round,line cap=round,font=\footnotesize,scale=.6]
			\tikzset{declare function={
					x1=-3.5;x2=4.5;y1=-1.5;y2=6.5;% Miền hình vẽ
					a=2;b=-1; m=1;n=-1;
					f(\x)=(a*(\x)+b)/(m*(\x)+n);	
					a1=a/m;b1=(b-a*n/m)/m;
					g(\x)=a/m;
					q1=-(n/m)-.1;q2=-n/m+.1;
			}}
			\begin{scope}
				\clip (x1,y1)rectangle(x2,y2);
				\draw [smooth,samples=300,domain=x1:q1]plot(\x,{f(\x)});
				\draw [smooth,samples=300,domain=q2:x2]plot(\x,{f(\x)});
				\draw [smooth,samples=100,domain=x1:x2]plot(\x,{g(\x)});
				\draw (-n/m,y1)--(-n/m,y2);
			\end{scope}
			\draw[->] (0,y1)--(0,y2)node[left]{$y$};
			\draw[->] (x1,0)--(x2,0)node[above]{$x$};
			\filldraw(0,0)circle(1.2pt)node[below left]{$O$};
			\draw[dashed]
			;
			\foreach \x/\goc in {1/-60}{
				\draw[fill=black] (\x,0)circle(1.2pt) node[shift={(\goc:2.8mm)}]{$\x$};
			}
			\foreach \y/\goc in {1/180,2/150}{
				\draw[fill=black] (0,\y)circle(1.2pt)node[shift={(\goc:2.8mm)}]{$\y$};
			}
		\end{tikzpicture}
	}
	\loigiai{
		\immini {
			Đồ thị hàm số $ y=\dfrac{2\left|x\right|-1}{\left|x\right|-1}$ có được bằng cách lấy đối xứng phần đồ thị phía bên phải trục tung qua trục $Oy$ và xóa phần đồ thị phía bên trái trục tung của đồ thị hàm số $ y=\dfrac{2x-1}{x-1}$.\\
			Dựa vào đồ thị ta thấy hàm số $ y=|x|^3-3|x|+1 $ đồng biến trên khoảng $(-\infty;-1)$ và $(-1;0)$.
		}
		{\begin{tikzpicture}[>=stealth,line join=round,line cap=round,font=\footnotesize,scale=.6]
				\tikzset{declare function={
						x1=-3.5;x2=4.5;y1=-1.5;y2=6.5;% Miền hình vẽ
						a=2;b=-1; m=1;n=-1;
						f(\x)=(a*(\x)+b)/(m*(\x)+n);	
						a1=a/m;b1=(b-a*n/m)/m;
						g(\x)=a/m;
						q1=-(n/m)-.1;q2=-n/m+.1;
				}}
				\begin{scope}
					\clip (x1,y1)rectangle(x2,y2);
					\draw [smooth,samples=300,domain=x1:x2]plot(\x,{f(abs(\x))});
					\draw [smooth,samples=300,domain=q2:x2]plot(\x,{f(abs(\x))});
					\draw [smooth,samples=100,domain=x1:x2]plot(\x,{g(\x)});
					\draw (-n/m,y1)--(-n/m,y2);
				\end{scope}
				\draw[->] (0,y1)--(0,y2)node[left]{$y$};
				\draw[->] (x1,0)--(x2,0)node[above]{$x$};
				\filldraw(0,0)circle(1.2pt)node[below left]{$O$};
				\draw[dashed]
				;
				\foreach \x/\goc in {1/-60,-1/-130}{
					\draw[fill=black] (\x,0)circle(1.2pt) node[shift={(\goc:2.8mm)}]{$\x$};
				}
				\foreach \y/\goc in {1/180,2/150}{
					\draw[fill=black] (0,\y)circle(1.2pt)node[shift={(\goc:2.8mm)}]{$\y$};
				}
		\end{tikzpicture}}
	}
\end{bt}

% nguồn: 2026-11-16_12_2D1_TL_Ham_so (nhập 16/11/2026)
\begin{bt}%[2D1V5-2]%[ma:cd4b0cb7]
% nguon: Archive/2024/TaiLieuGiangDay/1.HamSo/1. K12 - 4in1 - Hamso/ThubaiGoc/Bai5-Dang5.tex #2317

	\immini{
		Cho đồ thị của hàm số $y=f(x)=\log_4 x$ có đồ thị như hình vẽ. Từ đó, hãy suy ra đồ thị của hàm số $y=\left|f(x-2)\right|$.
	}{
		\begin{tikzpicture}[>=stealth,line join=round,line cap=round,font=\footnotesize,scale=.8]
			\tikzset{declare function={
					x1=-2.2;x2=4.5;y1=-2.5;y2=2.5;
					a=4;b=0;
					f(\x)=ln(\x+b)/ln(a);
					x0=-b+.02;
			}}
			\begin{scope}
				\clip (x1,y1)rectangle(x2,y2);
				\draw [smooth,samples=200,domain=x0:x2]plot(\x,{f(\x)});
			\end{scope}
			\draw[->] (0,y1)--(0,y2)node[left]{$y$};
			\draw[->] (x1,0)--(x2,0)node[above]{$x$};
			\filldraw(0,0)circle(1.2pt)node[below left]{$O$};
			\foreach \x/\goc in {1/90}{
				\draw[fill=black] (\x,0)circle(1.2pt) node[shift={(\goc:2.8mm)}]{$\x$};
			}
			\foreach \y/\goc in {}{
				\draw[fill=black] (0,\y)circle(1.2pt)node[shift={(\goc:2.8mm)}]{$\y$};
			}
		\end{tikzpicture}	
	}	
	\loigiai{	
		Trước tiên ta vẽ đồ thị hàm số $y=\log_4 x$.\\	
		Ta thấy đồ thị hàm số $y=|\log_4x|$ có được bằng cách:
		\begin{itemize}
			\item Giữ nguyên phần đồ thị của hàm số $y=\log_4x$ nằm trên trục hoành.
			\item Lấy đối xứng qua trục hoành phần đồ thị của $y=\log_4x$ nằm dưới trục hoành.
		\end{itemize}
		Do đó đồ thị hàm số $y=|\log_4x|$ là
		\begin{center}
			\begin{tikzpicture}[>=stealth,line join=round,line cap=round,font=\footnotesize,scale=.8]
				\tikzset{declare function={
						x1=-2.2;x2=4.5;y1=-2.5;y2=2.5;
						a=4;b=0;
						f(\x)=ln(\x+b)/ln(a);
						x0=-b+.02;
				}}
				\begin{scope}
					\clip (x1,y1)rectangle(x2,y2);
					\draw [smooth,samples=200,domain=x0:x2]plot(\x,{abs(f(\x))});
				\end{scope}
				\draw[->] (0,y1)--(0,y2)node[left]{$y$};
				\draw[->] (x1,0)--(x2,0)node[above]{$x$};
				\filldraw(0,0)circle(1.2pt)node[below left]{$O$};
				\foreach \x/\goc in {1/90}{
					\draw[fill=black] (\x,0)circle(1.2pt) node[shift={(\goc:2.8mm)}]{$\x$};
				}
				\foreach \y/\goc in {}{
					\draw[fill=black] (0,\y)circle(1.2pt)node[shift={(\goc:2.8mm)}]{$\y$};
				}
			\end{tikzpicture}
		\end{center}
		Đồ thị hàm số $y=\left|\log_4(x-2)\right|$ có được bằng cách tịnh tiến đồ thị hàm số $y=\left|\log_4x\right|$ sang phải $2$ đơn vị theo phương của trục $Ox$.
		\begin{center}
			\begin{tikzpicture}[>=stealth,line join=round,line cap=round,font=\footnotesize,scale=.8]
				\tikzset{declare function={
						x1=-2.2;x2=4.5;y1=-2.5;y2=2.5;
						a=4;b=-2;
						f(\x)=ln(\x+b)/ln(a);
						x0=-b+.02;
				}}
				\begin{scope}
					\clip (x1,y1)rectangle(x2,y2);
					\draw [smooth,samples=200,domain=x0:x2]plot(\x,{abs(f(\x))});
				\end{scope}
				\draw[->] (0,y1)--(0,y2)node[left]{$y$};
				\draw[->] (x1,0)--(x2,0)node[above]{$x$};
				\filldraw(0,0)circle(1.2pt)node[below left]{$O$};
				\draw (2,y1)--(2,y2);
				\foreach \x/\goc in {2/-110,3/-90}{
					\draw[fill=black] (\x,0)circle(1.2pt) node[shift={(\goc:2.8mm)}]{$\x$};
				}
				\foreach \y/\goc in {}{
					\draw[fill=black] (0,\y)circle(1.2pt)node[shift={(\goc:2.8mm)}]{$\y$};
				}
			\end{tikzpicture}
		\end{center}	
	}
\end{bt}

% nguồn: chuyen_dang (nhập 22/11/2026)
\begin{ex}%[2D1V5-2]%[ma:60b98637]
	 Cho đồ thị hàm số $y=f(x)$ dưới đây. Xét tính đúng sai của các khẳng định sau
	\begin{center}
		\begin{tikzpicture}[line join=round,line cap=round,>=stealth,xscale=1.2,yscale=1]
				 \tikzset{every node/.style={scale=0.7}}
				\def\a{-2.5}  \def\c{3.5}
				 \def\b{-1} \def\d{4}
				\def\f{3^(\x)} 
				\begin{scope}
					\clip (\a,\b)  rectangle (\c,\d);
					\draw[samples=200,domain=\a:\c,smooth,variable=\x] plot  (\x,{\f});
				\end{scope}
				\draw[->] (\a,0)--(\c+0.5,0) node[below] {$x$}; 
				\draw[->] (0,\b)--(0,\d+0.5) node[left]  {$y$};
				\draw (0,0) node  [below left] {$O$};
				 \foreach \x/\g in {1/60}
				 \fill[black] (\x,0) circle (1pt) + (\g:2.5mm) node {$\x$};
				\foreach \y/\h in {1/180,3/180} 
				\fill[black]  (0,\y) circle (1pt)+ (\h:2.5mm) node {$\y$};
				\foreach \x/\y  in {1/3}
				\draw[dashed,thin]  (\x,0)--(\x,\y)--(0,\y);
				 \node at (1.8,3.8) {$y=a^x$};
			\end{tikzpicture}
	\end{center}
	\choiceTF 
		{Hàm số đã  cho nghịch biến trên khoảng $(0;1)$}
		 {\True Hàm số cho bởi công thức $y=3^x$}
		{\True Đồ thị hàm số đã cho cắt đường thẳng $y=\dfrac{1}{3}$  tại điểm có hoành độ không âm}
		{Đồ thị hàm số đã cho cắt đường thẳng $y=-x+1$ tại điểm  có hoành độ dương}
	\loigiai{ 
		\begin{itemchoice}
			\itemch Dựa vào đồ thị hàm số ta  thấy hàm số đã cho đồng biến trên $\mathbb{R}$.
			 \itemch Đồ thị hàm số đi qua điểm $(1;3)$ suy ra $a=3$.
			\itemch Phương trình hoành độ giao  điểm của $y=3^x$ và đường thẳng $y=\dfrac{1}{3}$ là
			$$3^x=\dfrac{1}{3}\Leftrightarrow x=-1<0.$$ 
			\itemch Phương trình hoành độ giao điểm của $y=3^x$ và đường thẳng  $y=-x+1$ là $$3^x=-x+1\quad (1)$$
			Ta có hàm số $y=3^x$ đồng biến  trên $\mathbb{R}$ và $3^x>0$,$\forall x\in\mathbb{R}$.\\
			Đường thẳng $y=-x+1$ có hệ  số $a=-1<0$ nên nghịch biến trên $\mathbb{R}$.\\
			Ta lại có $3^0=-0+1$ nên phương trình (1) có nghiệm duy nhất  $x=0$.
		\end{itemchoice}
	 }
\end{ex}

% nguồn: chuyen_dang (nhập 22/11/2026)
\begin{bt}%[2D1H5-2]%[Huỳnh Đức Vũ]%[0D2K1]%[ma:dc2201fd]
	\immini{Cho hàm số $y=f(x)$ có đồ thị như hình vẽ bên.
		\begin{enumerate}[a)]
			\item Tìm các khoảng đồng biến, các khoảng nghịch biến của hàm số.
			\item Tìm giá trị lớn nhất, giá trị nhỏ nhất của hàm số trên đoạn $\left[0; 2\right]$.
	\end{enumerate}}{
		\begin{tikzpicture}[>=stealth,line join=round,line cap=round,font=\footnotesize,scale=0.6]			
			\draw[->,color=black] (-2.02,0.) -- (2.34,0.);
			\foreach \x in {-1,2}
			\draw[shift={(\x,0)},color=black] (0pt,2pt) -- (0pt,-2pt) node[below] {$\x$};
			\foreach \x in {-2,1}
			\draw[shift={(\x,0)},color=black] (0pt,2pt) -- (0pt,-2pt) node[above] {$\x$};
			\draw[color=black] (2.02,0.08) node [anchor=south west] {$x$};
			\draw[->,color=black] (0.,-2.68) -- (0.,2.6);
			\foreach \y in {-2}
			\draw[shift={(0,\y)},color=black] (2pt,0pt) -- (-2pt,0pt) node[left] {$\y$};
			\foreach \y in {2}
			\draw[shift={(0,\y)},color=black] (2pt,0pt) -- (-2pt,0pt) node[above left] {$\y$};
			\draw[color=black] (0.1,2.6) node [anchor=west] { $y$};
			\draw[color=black] (0pt,-10pt) node[right] {$\tiny  {\text{O}}$};
			\clip(-2.02,-2.68) rectangle (2.34,2.66);
			\draw[smooth][line width=0.8pt,color=black,smooth,samples=100,domain=-2.02:2.02] plot(\x,{(\x^3.0-3.0*\x});
			\draw [dashed] (-1.,2.)-- (0,2)--(2,2)--(2,0);
			\draw [dashed] (-1.,2.)-- (-1.,0.);
			\draw [dashed] (1.,0.)-- (1.,-2.);
			\draw [dashed] (1.,-2.)-- (0.,-2.);
			\fill (1,-2) circle(1pt) 
			(-1,2) circle(1pt)(2,2) circle(1pt);	
	\end{tikzpicture}}
	\loigiai{ \begin{enumerate}[a)]
			\item
			Từ đồ thị, ta có hàm số đồng biến trên các khoảng $(-\infty; -1)$ và $(1; +\infty)$, nghịch biến trên khoảng $(-1;1)$.
			\item Từ đồ thị ta thấy trên đoạn $[0;2]$ hàm số đã cho có giá trị nhỏ nhất bằng $-2$ và giá trị nhỏ nhất bằng $2$.
	\end{enumerate}}
\end{bt}
